\titleformat{name=\chapter,numberless}[display]
{\normalfont\huge\bfseries}
{}
{1em}
{}

\titleformat{\chapter}[display]
{\normalfont\huge\bfseries}
{Book \Ordinalstring{chapter}}
{1em}
{}


\titleformat{name=\section,numberless}
{\normalfont\Large\bfseries}	 	% 1. Format (shape/font)
{}                          	 	% 2. Label (empty if you want no number)
{1em}                       	 	% 3. Sep (vertical space in 'display' mode)
{}


\titleformat{\section}[display]
{\normalfont\Large\bfseries}	 	% 1. Format (shape/font)
{the \Ordinalstring{section} Chapter}      % 2. Label (empty if you want no number)
{1em}                       	 	% 3. Sep (vertical space in 'display' mode)
{} 	% 4. Code before title body



\titleformat{\subsection}[runin]
{\itshape\normalsize}	 	% 1. Format (shape/font)
{\arabic{subsection}}
{1em}
{\sidenote} 	% 4. Code before title body

\titleformat{\subsection}[runin]
{\itshape\normalsize}	 	% 1. Format (shape/font)
{\arabic{subsection}}
{1em}
{\sidenote\} 	% 4. Code before title body



%\titleformat{\subsection}[display]
%{\normalfont\Large\bfseries}	 	% 1. Format (shape/font)
%{}                          	 	% 2. Label (empty if you want no number)
%{1em}                       	 	% 3. Sep (vertical space in 'display' mode)
%{} 	% 4. Code before title body

\ExplSyntaxOn
% Allocate internal token storage strings
\tl_new:N \g_saved_subsec_title_tl


% Allocate internal token storage strings
\tl_new:N \gSavedSubsecLongTl
\tl_new:N \gSavedSubsecShortTl

% 1. Redefine \subsection to accept \subsection[Short]{Long}
\let\origsubsection\subsection
\RenewDocumentCommand{\subsection}{ o m }{
	\tl_gset:Nn \gSavedSubsecLongTl {#2}
	\IfValueTF{#1}
	{ \tl_gset:Nn \gSavedSubsecShortTl {#1} }
	{ \tl_gset:Nn \gSavedSubsecShortTl {#2} }

	\peek_meaning_ignore_spaces:NTF \subsubsection
	{ \catSkipToSubsub:w }
	{
		\IfValueTF{#1}
		{ \origsubsection[#1]{#2} }
		{ \origsubsection{#2} }
	}
}

% Cleanly consume the \subsubsection token
\cs_new:Npn \catSkipToSubsub:w \subsubsection #1 {
	\catParseSubsub:w #1
}

% Helper to check for subsubsection optional argument
\NewDocumentCommand{\catParseSubsub}{ o m }{
	\IfValueTF{#1}
	{ \catExecuteMerge:nnnn {\gSavedSubsecShortTl} {\gSavedSubsecLongTl} {#1} {#2} }
	{ \catExecuteMerge:nnnn {\gSavedSubsecShortTl} {\gSavedSubsecLongTl} {#2} {#2} }
}

% Process counters, ToC entries, and push heading to the margin
% #1 = Subsec Short, #2 = Subsec Long, #3 = Subsub Short, #4 = Subsub Long
\cs_new:Npn \catExecuteMerge:nnnn #1 #2 #3 #4 {
\refstepcounter{subsection}
\refstepcounter{subsubsection}

% Track cleanly in the Table of Contents
\addcontentsline{toc}{subsection}{\protect\numberline{\thesubsection}#1}
\addcontentsline{toc}{subsubsection}{\protect\numberline{\thesubsubsection}#3}

% 1. Place the reference anchor mark in the restricted heading zone
\sidenotemark

% 2. Safely defer the text placement until LaTeX drops back to outer paragraph mode
\bgroup
\toks0={
\sidenotetext{
	\raggedleft
	\normalfont\itshape\bfseries\small
	\thesubsubsection \quad #2 \\ \vspace{0.2ex} \rule{1cm}{0.4pt} \\ \vspace{0.2ex} #4
}
}
\expandafter\egroup\the\toks0

\ignorespaces
}
\ExplSyntaxOff
